To probe violation of conditional independence, we add a traceless correlation perturbation that preserves the local conditional marginals while increasing $\nu=I(E_i:E_j|X)$.  Figure~\ref{fig:hardwareinformed}(b) shows that the population decoder-direction error is approximately linear in $\sqrt{\nu}$ over the tested family,
\[
\sin\theta_{\rm pop}\simeq 1.66\sqrt{\nu}+0.014,
\qquad R^2=0.995,
\]
with $\nu$ measured in bits.  Finite-shot estimates track the same population degradation.  This is consistent with the square-root perturbation scale entering the Pinsker--Wedin robustness argument, while also showing that the analytical worst-case bound is conservative.

\subsection{Multiclass ambiguity and third-view finite-shot recovery}

The commuting three-class counterexample of Appendix~\ref{app:counterexample} reconstructs the same two-view distribution under two inequivalent latent decompositions to $1.39\times10^{-17}$ numerical precision, yet the corresponding MAP decoders differ.  This ambiguity therefore survives infinite data.  We then append a well-conditioned third stochastic view and estimate the resulting three-view tensor from finite samples.  As shown in Fig.~\ref{fig:hardwareinformed}(c), exact decoder recovery rises from $20\%$ at $10^3$ samples to $92.5\%$ at $3\times10^4$ and $100\%$ at $10^5$, while the spectral failure rate falls to zero.  The latent components are reconstructed without access to the true factors; a common permutation is matched to the planted model only after reconstruction for error reporting, as required by the identifiability statement.

\begin{figure*}[t]
\centering
\includegraphics[width=0.99\textwidth]{fig9_hardware_informed_identifiability.pdf}
\caption{Hardware-informed numerical checks matched to the identifiability theory.  (a) Median binary decoder sin-angle error versus the number of physical events for three contrast amplitudes.  The observed slopes are close to $N^{-1/2}$; the inset shows the approximate collapse of $c^2\sqrt N$ times the median error.  (b) Controlled violation of conditional independence.  Population and finite-shot decoder errors grow approximately linearly with $\sqrt{I(E_i:E_j|X)}$ over the tested family; only the portion of the conservative Pinsker--Wedin bound below unity is displayed.  (c) Multiclass transition.  Two views admit an exact population ambiguity with inequivalent decoders, whereas a well-conditioned third view yields increasingly reliable finite-sample recovery.  These panels test scaling and identifiability transitions, not quantum advantage or hardware calibration.}
\label{fig:hardwareinformed}
\end{figure*}

\subsection{Near-SBS spectral anchor}

The three-class test shown in Fig.~\ref{fig:nearsbs} uses an ideal orthogonal $B$ record and three distinct conditional qubit states on $A$.  A single probe has expectations $(0.56,0.14,-0.42)$, exactly reproduced as the eigenvalues of the whitened operator.  Mixing in a generic bipartite perturbation yields projector errors from $5.7\times10^{-5}$ at perturbation $10^{-4}$ to $1.8\times10^{-2}$ at perturbation $3\times10^{-2}$.

\subsection{Microscopic collision-model simulation}

We next generate the environmental records from sequential system--environment dynamics rather than inserting the conditional record states directly.  The system starts in $|+\rangle_S$ and five environmental qubits in $|0\rangle^{\otimes5}$.  Each fragment interacts through a controlled rotation
\[
U_{SE_j}(\chi)=|0\rangle\!\langle0|_S\otimes I_j
+|1\rangle\!\langle1|_S\otimes R_y^{(j)}(\chi),
\]
after which every fragment is independently conjugated by an unknown local $SU(2)$ rotation.  The system is traced out before measurement.  The learner receives only finite-shot environmental Pauli outcomes from a 27-setting orthogonal-array schedule, with independent symmetric readout flips of probability $q=0.02$.  The local branch distinguishability generated by the dynamics is $D=\sin(\chi/2)$.

Figure~\ref{fig:collisionphase} uses the strict success criterion $\max_j\theta_j\le5^\circ$: every learned fragment axis must lie within $5^\circ$ of its hidden decoder.  The resulting phase diagram separates weak records, for which additional shots cannot compensate for poor conditioning on the tested budget, from strong records that become reliably identifiable as the number of shots grows.  For example, at $D=0.707$ the all-fragment success probability increases from $0.14$ at 128 shots per setting to $0.83$ at 512, while at $D=1$ it rises from $0.31$ at 64 shots to $1.00$ at 512.  This is a finite-shot virtual experiment under a simple measurement-noise model, not a calibrated device simulation.

\begin{figure}[t]
\centering
\includegraphics[width=\columnwidth]{fig10_collision_recovery_phase_diagram.pdf}
\caption{Microscopic system--environment collision-model simulation.  Entries show the probability that \emph{all five} learned decoder axes satisfy $\theta_j\le5^\circ$, using 27 environment-only Pauli settings and $2\%$ symmetric readout flips.  The collision angles corresponding to $D=0.174,0.342,0.500,0.707,0.866,1.000$ are $20^\circ,40^\circ,60^\circ,90^\circ,120^\circ,180^\circ$, respectively.  Stronger records and larger shot budgets move the estimator into a stable recovery regime.}
\label{fig:collisionphase}
\end{figure}

\subsection{Chen-derived hidden-readout benchmark}

To retain contact with an external Darwinism architecture, we include the six-photon setup of Chen \emph{et al.} as a simulation benchmark rather than an experimental reanalysis \cite{chen2019}.  Each environmental photon is independently conjugated by a hidden local $SU(2)$ rotation.  The learner receives no system labels or hidden bases.  In the ideal all-strong-record setting, the earlier agreement-based search reaches unit output richness with residual worst-pair disagreement $0.00193\pm0.00093$ across 12 randomized hidden-basis worlds.  In the mixed setting with three strong and two weak records, the learned disagreement $0.27689\pm0.00496$ lies close to the physical oracle floor $0.27487$; the best three-fragment subset is the three strong records in all 12 population worlds.

\begin{figure*}[t]
\centering
\begin{minipage}{0.43\textwidth}
\centering
\begin{tikzpicture}[font=\sffamily\small,>=Latex,
 ph/.style={circle,draw=envteal,fill=envteal!8,minimum size=7.5mm},
 hid/.style={rounded corners=1.5pt,draw=accentorange,fill=accentorange!8,minimum width=12mm,minimum height=6mm,align=center}]
\node[font=\bfseries] at (2.9,2.35) {(a) Chen-derived SIMULATION benchmark};
\node[circle,draw=sysblue,very thick,fill=sysblue!8,minimum size=10mm] (S) at (.6,.5) {$S$};
\node[ph] (E2) at (1.8,1.55) {$E_2$};
\node[ph] (E3) at (2.7,1.05) {$E_3$};
\node[ph] (E4) at (3.2,.2) {$E_4$};
\node[ph] (E5) at (2.7,-.65) {$E_5$};
\node[ph] (E6) at (1.8,-1.15) {$E_6$};
\foreach \E in {E2,E3,E4,E5,E6}{\draw[->,semithick,draw=envteal] (S)--(\E);}
\node[hid] (u2) at (4.55,1.55) {$U_2$};
\node[hid] (u3) at (4.55,.90) {$U_3$};
\node[hid] (u4) at (4.55,.25) {$U_4$};
\node[hid] (u5) at (4.55,-.40) {$U_5$};
\node[hid] (u6) at (4.55,-1.05) {$U_6$};
\draw[->] (E2)--(u2); \draw[->] (E3)--(u3); \draw[->] (E4)--(u4); \draw[->] (E5)--(u5); \draw[->] (E6)--(u6);
\node[font=\scriptsize,align=center,text=deepgray] at (2.8,-1.85) {independent hidden local $SU(2)$ bases;\\no access to $S$ or the hidden rotations};
\end{tikzpicture}
\end{minipage}\hfill
\begin{minipage}{0.54\textwidth}
\centering
{\bfseries\color{accentred} SIMULATION}\par\vspace{0.4mm}
\includegraphics[width=\linewidth]{fig7_simulation_benchmark.pdf}
\end{minipage}
\caption{External hidden-readout simulation stress test.  (a) The record architecture follows the six-photon Quantum-Darwinism experiment of Chen \emph{et al.}, but each environmental photon is independently conjugated by an unknown local $SU(2)$ rotation before readout.  (b) In simulation, the learned decoder remains close to the physical oracle floor in both the all-strong-record and mixed strong/weak-record settings.  This is a benchmark derived from the published architecture, not a reanalysis of raw experimental events.}
\label{fig:chen}
\end{figure*}

\subsection{Operational resource tradeoff}

